% Options for packages loaded elsewhere
\PassOptionsToPackage{unicode}{hyperref}
\PassOptionsToPackage{hyphens}{url}
%
\documentclass[
  11pt,
]{article}
\usepackage{amsmath,amssymb}
\usepackage{iftex}
\ifPDFTeX
  \usepackage[T1]{fontenc}
  \usepackage[utf8]{inputenc}
  \usepackage{textcomp} % provide euro and other symbols
\else % if luatex or xetex
  \usepackage{unicode-math} % this also loads fontspec
  \defaultfontfeatures{Scale=MatchLowercase}
  \defaultfontfeatures[\rmfamily]{Ligatures=TeX,Scale=1}
\fi
\usepackage{lmodern}
\ifPDFTeX\else
  % xetex/luatex font selection
\fi
% Use upquote if available, for straight quotes in verbatim environments
\IfFileExists{upquote.sty}{\usepackage{upquote}}{}
\IfFileExists{microtype.sty}{% use microtype if available
  \usepackage[]{microtype}
  \UseMicrotypeSet[protrusion]{basicmath} % disable protrusion for tt fonts
}{}
\makeatletter
\@ifundefined{KOMAClassName}{% if non-KOMA class
  \IfFileExists{parskip.sty}{%
    \usepackage{parskip}
  }{% else
    \setlength{\parindent}{0pt}
    \setlength{\parskip}{6pt plus 2pt minus 1pt}}
}{% if KOMA class
  \KOMAoptions{parskip=half}}
\makeatother
\usepackage{xcolor}
\usepackage[margin=1in]{geometry}
\setlength{\emergencystretch}{3em} % prevent overfull lines
\providecommand{\tightlist}{%
  \setlength{\itemsep}{0pt}\setlength{\parskip}{0pt}}
\setcounter{secnumdepth}{-\maxdimen} % remove section numbering
\ifLuaTeX
  \usepackage{selnolig}  % disable illegal ligatures
\fi
\usepackage{bookmark}
\IfFileExists{xurl.sty}{\usepackage{xurl}}{} % add URL line breaks if available
\urlstyle{same}
\hypersetup{
  pdftitle={The universal contact property for products of finite rank-one balanced graphs},
  pdfauthor={Xean research manuscript},
  hidelinks,
  pdfcreator={LaTeX via pandoc}}

\title{The universal contact property for products of finite rank-one
balanced graphs}
\author{Xean research manuscript}
\date{4 October 2026}

\begin{document}
\maketitle

\subsection{Statement, coefficient convention, and residual
scope}\label{statement-coefficient-convention-and-residual-scope}

Let \(G_i\subset N_{i,\mathbb{R}}\), \(i=1,2\), be nonempty finite
embedded rational pure graphs. Their edge weights \(w_{i,e}\) are
positive real numbers. At every vertex \(v\), including inserted
bivalent vertices, assume \[
 K_{i,v}=\{c\geq0:\sum_e c_e v_e=0\}
       =\mathbb{R}_{\geq0}(w_{i,e})_e,                         \tag{1}
\] where \(v_e\) are the actual saturated primitive outgoing directions.
The graph structure can have bounded edges, arbitrary cycles, full
affine lines and unbounded ends, and finitely many connected components.
An incident direction means an actual carrier germ; a \textbf{branched
junction} has at least three such directions. Give \(X=G_1\times G_2\)
the lattice \(N_1\oplus N_2\) and actual product weights.

\textbf{Universal graph-product theorem.} The following property
holds if and only if at least one factor has no branched junction:

For \textbf{every} continuous real-PA actual-PC reference \(\gamma\)
positive at infinity and \textbf{every} bounded continuous real-PA
obstacle \(u\), the actual raw envelope \(P_\gamma(u)\) is PAPC on a
finite rational spatial structure, is uniformly decreasingly
regularizable in the two-sided \(\gamma+O(1)\) class, and its complete
recursive operator satisfies \[
 \int_X(P_\gamma(u)-\gamma-u)\,
                \mathrm{MA}_{\mathrm{poly},\mathbb{R}}(P_\gamma(u))=0.
                                                               \tag{2}
\] Real PA means real affine coefficients on \textbf{finite rational
spatial domains}, not arbitrary real walls. Raw PC and PAPC have all the
tests specified below. The operator is the full two-stage intersection
operator, with no division by \(2!\).

In the negative direction, choose any branched vertex \(A\) of \(G_1\)
and \(B\) of \(G_2\). Write their primitive positive integer relations
as \(\lambda_i\) and \(\mu_j\), and their actual weights as
\(a\lambda_i\), \(b\mu_j\), \(a,b>0\). The proof constructs permitted
\(\gamma,u\) on the \textbf{entire original product} with \[
 P_\gamma(u)=\gamma\in\mathrm{PAPC}_{\mathbb{R}}
              \cap\mathrm{PC}_{\mathrm{reg},\mathbb{R}},\qquad
 \int_X(P_\gamma(u)-\gamma-u)\,\mathrm{MA}_{\mathrm{poly},\mathbb{R}}
                  (P_\gamma(u))=-8ab<0.                       \tag{3}
\] This is an existence of a bad reference on every two-branched-factor
product, not a classification of its individual references.

Original P05.2 asks for a characterization of fixed arbitrary balanced
pairs \((X,\gamma)\). That atom remains open. The theorem only
characterizes the displayed universal property in this defined
graph-product class. It neither settles references individually on two
branched factors nor the remaining higher-dimensional singular
contact/strict-positivity question, literal coefficients, analytic
realization, extra affine factors, or P06.2. This is a reviewed partial theorem;
predecessor notes, reviews and executions are not theorem premises.

\subsection{Definitions, source checks, and actual intersection
algebra}\label{definitions-source-checks-and-actual-intersection-algebra}

The controlling book files are Chapter 72,
\texttt{ch72\_major\_open\_problems.tex} (polyhedral section, original
problem and residual scope), and Appendix M,
\texttt{appM\_ch18\_frontiers.tex} (analytic zero-current catalogue, no
product theorem used). Their complete paths and fresh snapshots are in
\texttt{source-hashes.json}. BMPS, \emph{Monge--Ampère measures on
balanced polyhedral spaces}, arXiv:2603.08664v1, dated 10 March 2026,
supplies definitions only. The explicit extension used here has real
function coefficients and real cycle weights; BMPS Definition 2.28
itself has integral slopes/rational constants. On any fixed finite
rational complex, affine-piece continuity is a rational linear system;
its real solution space has a rational basis, each made literal by
clearing slope denominators. Thus this convention is the real scalar
extension on that complex.

A PCC (BMPS Definitions 3.2--3.3) has rational endpoints
\(x_i\in X\cap N_{\mathbb{Q}}\), real coefficients \(\alpha_i\geq0\),
\(\sum\alpha_i=1\), and the actual, possibly irrational, centre
\(x=\sum\alpha_i x_i\). On some allowed rational structure there must be
faces \(\tau\preceq\sigma_i\) containing \(x\) and \(x_i\). A finite raw
PC function satisfies Jensen for \textbf{every} such combination, across
all allowed structures and remote endpoints. No continuity of raw
competitors is assumed.

With \(f=\gamma+u\), the raw envelope (Definition 6.2) is \[
 P_\gamma(u)(x)=\sup\{q(x):q\in\mathrm{PC}(X),\ q\leq f,
                          \ q\leq\gamma+C_q\text{ for some constant}\}.
                                                               \tag{4}
\] The last condition is only upper growth. The regularization class has
two-sided growth (Definitions 3.14, 3.23). Positive at infinity
(Definition 4.15) means strictly positive slope on every nonzero
recession vector of every affine reference cell. PAPC (Definition 3.1,
here with real weights) tests \(h\cdot z\geq0\) for every nonnegative
balanced real local cycle on every ordinary open, in every positive
dimension. On this two-dimensional carrier only dimensions one and two
can be nonzero.

Use the actual lattice
\(N_\sigma=N\cap\operatorname{span}_{\mathbb{R}}(\sigma-\sigma)\). For a
facet \(\tau\prec\sigma\), \(n_{\sigma/\tau}\) lifts the primitive
positive class in \(N_\sigma/N_\tau\). The divisor formula (Definition
2.31) is \[
 (h\cdot z)(\tau)=\sum_{\sigma\succ\tau}z(\sigma)h_\sigma(n_{\sigma/\tau})
       -h_\tau\left(\sum_{\sigma\succ\tau}z(\sigma)n_{\sigma/\tau}\right).
                                                               \tag{5}
\] The second sum is tangential by balancing. Changes of lift cancel
because the linear parts agree on tangents. Formula (5) is real
bilinear, depends only on incident germs, and annihilates ambient affine
functions.

Here is a direct reconstruction of refinement and balancing, rather than
an appeal to source intersection theory. For a new tested facet in a
refinement, take its minimal old face. If its dimension exceeds the
cycle dimension, inherited incident cycle weights are zero. If its
dimension equals the cycle dimension, the two inherited pieces in that
plane have equal weights and affine parts and opposite primitive
quotient normals, so cancel. If its dimension equals the tested facet
dimension, the actual tangent lattice and incident normal classes agree
with the old ones, so the old coefficient is inherited. These cases
prove refinement compatibility and inherited balancing. They apply
successively to both divisors; in particular new vertices in affine
edges cannot create a zero-cycle coefficient.

For balancing of the first divisor of a two-cycle, translate a vertex to
zero and take a pointed rational local fan, adding harmless rational
cuts if necessary. In a two-cone with primitive rays \(r,s\), let \[
 \chi_{rs}=[N_C:\mathbb Zr+\mathbb Zs],\qquad \omega_{rs}=z(C).
\] Extending primitive \(r\) to a basis of the saturated plane lattice
shows that its inward primitive quotient normal is represented by
\(s/\chi_{rs}\). A nonintegral representative has only a tangential
discrepancy from an integral lift, which cancels in (5). Top balancing
and the first weight are \[
 \sum_{C\succ r}(\omega_{rs}/\chi_{rs})s=t_r r,
 \qquad D_h(r)=\sum_{C\succ r}(\omega_{rs}/\chi_{rs})h(s)-t_rh(r). \tag{6}
\] Summing \(D_h(r)r\), the first terms, once per unordered cone, are
\(\omega_{rs}(h(s)r+h(r)s)/\chi_{rs}\). Substituting top balancing into
the subtracted terms gives exactly the same sum. Hence
\(\sum_rD_h(r)r=0\). This proves first-divisor balancing for real
weights in every germ. The full second coefficient is \[
                       (h^2\cdot[X])\{0\}=\sum_rD_h(r)h(r).     \tag{7}
\] Every intermediate ray, including zero-weight or zero-value rays,
belongs to this formula. Definition 4.4's complete measure is the
resulting finite zero-cycle on structure vertices. On a unit complete
plane, \(|s|+|t|\) has first weight two on each axis ray, and (7) gives
eight: \(2!\) times the slope-square area four. No factorial division,
probability normalization, or singular-face Alexandrov substitution is
used.

If a continuous real PA function is actual PC, it is PAPC against
\textbf{all} required cycles. Indeed, refine a local cycle and the
function, choose a rational point \(x\) in the ordinary open and
relative interior of its tested facet \(\tau\), and first take
nonnegative rational balanced incident weights \(c_\sigma\). Put
\(S=\sum c_\sigma>0\) and
\(T=\sum c_\sigma n_{\sigma/\tau}\in N_{\tau,\mathbb{Q}}\). For small
rational \(\varepsilon>0\), \[
 x_\sigma=x+\varepsilon(n_{\sigma/\tau}-T/S)
\] lies in its actual incident cell and open: the correction is
tangential and the quotient direction is inward. These rational
endpoints, with coefficients \(c_\sigma/S\), give an actual PCC centred
at \(x\). Jensen gives
\(0\leq\varepsilon(\sum c_\sigma h_\sigma(n_{\sigma/\tau})-h_\tau(T))\),
precisely (5). The cone of nonnegative balanced incident vectors is a
rational polyhedral cone, so has rational generators, as proved below.
Linearity extends the result to every real incident vector. The zero
vector gives zero. This proves all ordinary-open, all-dimension PAPC
tests. Balancing just proved and this positivity allow both recursive
divisors.

For all envelopes constructed here, \(P+1/n\) is decreasing uniform PAPC
approximation in \(\gamma+O(1)\). Constants have zero divisor; the
approximant measures are identical to \(P^2\cdot[X]\). Thus the
regularized operator is concretely the already defined PA operator,
without an approximation, gluing, locality or continuity theorem
imported from the source. The full proof bodies of BMPS Lemmas
2.34--2.35, Lemma 3.4, Proposition 3.22 and Propositions 6.3 and 6.13
were inspected. Their external locality/intersection dependencies are
not used. There are \textbf{no essential external theorem premises} in
the argument that follows.

\subsection{Elementary finite-dimensional facts used in the
reconstruction}\label{elementary-finite-dimensional-facts-used-in-the-reconstruction}

A polyhedral cone has finitely many generators; for rational
inequalities these can be rational. Remove its lineality (a rational
subspace for rational data) and work in the pointed quotient. The sum of
its nonnegative defining forms is strictly positive away from zero:
otherwise all forms vanish on a nonzero vector and its negative. The
section of value one is bounded, since an unbounded sequence, normalized
to the unit sphere, would yield a nonzero cone vector of value zero. A
bounded polyhedron is the convex hull of its vertices: in the minimal
face of any nonvertex, a line ends on smaller faces, and induction
reduces dimension. Its vertices are the finitely many full-rank
active-system solutions, rational for rational data. Lift back and add
both signs of a lineality basis. Homogenizing a nonempty pointed
rational polyhedron therefore proves \[
 \sigma=\operatorname{conv}\{v_1,\ldots,v_l\}
                    +\operatorname{cone}\{r_1,\ldots,r_m\},     \tag{8}
\] with rational vertices and recession vectors and at least one vertex:
the height-one slice needs a positive-height extreme ray, whose
extremality gives an original vertex. The homogenized cone is pointed
because any line has height zero and would give original lineality. The
same proof holds over real data without rational coordinates.

Eliminating a variable from finitely many nonstrict linear inequalities
compares every lower bound with every upper bound and retains
variable-independent conditions. A one-sided family imposes no further
conditions. Repetition proves that projections of polyhedra are
polyhedra; variable coefficients stay rational if originally rational,
even with real constants. It also proves closedness of finite
convex-hull-plus-cone sets by projecting their coefficient
representation.

Needed separations can be proved by nearest points. A point outside a
nonempty closed convex set has a nearest point after restriction to a
sufficiently large compact ball. Differentiating squared distance along
segments gives a separating normal. A compact convex set disjoint from a
closed convex set has positive attained distance, with the same compact
restriction for the second set; differentiating in both sets strictly
separates them. If the latter is a cone, a finite supremum of the
separating functional on it is zero. This proves the cone/polytope
separations used here.

A finite rationally convex function on \(\mathbb{Q}\) has a unique
continuous convex extension on any interval interior. Rational endpoint
values bound it above on a rational compact interval. Reflection about a
rational centre bounds it below on a smaller interval. Bounds
\(L\leq F\leq M\) on a larger rational interval imply a Lipschitz bound
on a smaller one: extend the rational segment from \(y\) through \(z\)
by a rational factor \(R/|z-y|\) staying in the larger interval and use
Jensen, then reverse \(y,z\). Density gives the completion and
convexity. A finite convex function on \(\mathbb{R}\) is continuous by
this argument, and has a touching affine support at each point: its left
and right secant limits are finite, ordered, and any slope between them
gives the support.

For \(b(p)=\min_j(c_j+g_jp)\), the supergradient interval at \(p_0\) is
the convex hull of the active \(g_j\). Its directional derivative is the
minimum of their products with the direction. Separating a candidate
outside that interval disproves its supergradient property. Expanding
sums of finite minima proves that supergradients of nonnegative sums are
the corresponding weighted Minkowski sums, also when coefficients
vanish. The outward normal cone of a polytope is generated by active
inequality normals: an alleged normal outside their finitely generated
closed cone can be separated to give a direction satisfying the active
inequalities and increasing the alleged normal; a sufficiently small
step also satisfies inactive inequalities, a contradiction. These proofs
retain rational normal generators whenever inequality coefficients are
rational, even if constants are real.

\subsection{Graph geometry, real weights, and genuine
cofaces}\label{graph-geometry-real-weights-and-genuine-cofaces}

At a vertex the full kernel of the rational direction matrix is
one-dimensional. Otherwise perturb its strictly positive actual weight
vector by a sufficiently small nonproportional kernel vector; positivity
persists, contradicting (1). The kernel therefore has a primitive
positive integer generator \(\lambda_v\), and the actual incident vector
is \(a_v\lambda_v\), \(a_v>0\). At a bivalent vertex the directions are
opposite primitive vectors and weights equal. At a branched vertex there
is no opposite pair, since that pair would give another nonnegative
relation. No one-valent vertex can balance.

On a connected graph all edge weights are a single positive real scalar
times positive rational weights. Fix one edge weight. At each adjacent
vertex the ratios of incident weights equal ratios of positive integers
in \(\lambda_v\); propagate along paths through the finite graph.
Existing actual weights ensure consistency around every cycle. This
yields rational normalized weights \(r_e\) with \(w_e=c r_e\). Full-line
components have one arbitrary positive scalar weight. Rational
subdivision preserves these statements. This reduction is essential: it
makes the incidence elimination below have rational variable
coefficients, without supposing the given real scale rational.

Every connected unbranched component is a complete rational affine line.
Propagating through bivalent vertices keeps every edge on the same
affine line with equal weight. The component is a connected finite
closed union of intervals/rays on that line. Any finite endpoint of its
union would be one-valent, contrary to balancing. It therefore covers
the whole line; a bounded cycle cannot be embedded into this line as a
pure polyhedral graph. A component supplied already as a full line has
the same conclusion. Its tangent lattice is the saturated primitive line
lattice \(\mathbb Zv\), independently of the ambient lattice's
complementary directions.

For the all-coface PCC tests on a cylinder \(G\times\mathbb{R}\), any
convex cell projects to a convex subset of the graph. A segment cannot
pass a branched vertex since it would need an opposite pair there. Thus
its graph projection is in a single straight chain through bivalent
subdivisions or in a vertex. If the centre's graph coordinate is in a
straight-chain interior, all cofaces containing the common face remain
in that chain and the graph test is ordinary one-variable convexity. At
a branched centre \(v\), each coface projects into one incident straight
chain; every endpoint displacement has the form \(t_i v_e\),
\(t_i\geq0\), up to the next branched vertex. Grouping the barycentric
displacements gives \[
 b_e=\sum_{i\text{ in }e}\alpha_i t_i,
 \qquad\sum_e b_ev_e=0,\qquad b_e=c r_e,\ c\geq0.           \tag{9}
\] This argument uses actual cofaces on arbitrary rational structures,
not merely a prescribed product subdivision. Distant endpoints and
inserted bivalent vertices are included. Zero graph displacements cause
no contribution.

\subsection{Positive direction: full graph-cylinder reconstruction in
dimension
two}\label{positive-direction-full-graph-cylinder-reconstruction-in-dimension-two}

Suppose one factor is unbranched. Its components are complete rational
affine lines of constant positive real weights. Work first on one
connected component \(G\times\mathbb{R}_z\). Use its actual primitive
coordinate \(z\), and normalize graph weights to positive rational
\(r_e\); the actual product two-cycle is a positive scalar multiple of
these normalized weights. This scalar affects neither raw PC nor contact
and will be restored in the measure. No hypothesis of separability or
ambient convex extension of the data is made.

\textbf{Scalar graph lemma.} A finite function on rational graph points
satisfying rational-coefficient graph PCCs has a continuous
edgewise-convex completion \(k\) with finite outgoing derivatives and \[
                              \sum_e r_e k'_e(v+)\geq0.        \tag{10}
\] Complete it on edge interiors by the rational convex fact. The
one-sided limit at an endpoint is finite and at most its specified
endpoint value: endpoint secants give the upper bound and a fixed
interior secant gives a lower bound. For small rational \(t>0\), the
points \(v+t v_e\), with rational coefficients \(r_e/\sum r_e\), form an
actual vertex PCC. Its inequality and those endpoint upper bounds force
every limit to equal \(k(v)\). Each outgoing derivative is bounded above
by a finite secant. Dividing the same PCC by \(t\) prevents any
derivative from being \(-\infty\), and yields (10).

Conversely, any continuous edgewise-convex \(k\) satisfying (10), plus
\(pz\) for any real \(p\), satisfies \textbf{every} actual cylinder PCC.
On a chain this is convexity (the bivalent inequality joins the
derivatives). At a branched centre use (9) and
\(k(q_i)-k(v)\geq t_i k'_e(v+)\); summation gives
\(c\sum r_e k'_e(v+)\geq0\). Affine \(z\) terms cancel by the barycentre
identity. This includes arbitrary real centres/coefficients and every
coface structure.

For an arbitrary finite raw cylinder PC function \(h\), define on
rational graph points \[
 k_{h,p}(q)=\inf_{z\in\mathbb{Q}}(h(q,z)-pz).                 \tag{11}
\] It satisfies rational graph PCCs: choose rational \(z_i\) tending to
each endpoint infimum, apply the product PCC with rational coefficients
at \((q,\sum\alpha_i z_i)\), subtract its affine term, and take the
infimum. A \(-\infty\) endpoint is handled by varying that endpoint
alone. If (11) is finite at one rational graph point, it is finite
throughout the component. From an edge-interior finite centre any other
rational edge point can be a positively weighted endpoint, with the
other endpoint on the opposite side. From a finite vertex, an arbitrary
rational point \(v+t_ev_e\) on an incident edge is an endpoint in a
positive rational vertex PCC: choose positive rational distances on the
other edges and coefficients proportional to \(r_e/t_e\). No endpoint of
a finite centre's such inequality can be \(-\infty\). Propagation across
the connected graph includes every cycle and end. Thus the scalar lemma
applies whenever (11) is finite anywhere.

Put \(f=\gamma+u\), \(C=\|u\|_\infty\), and take a common finite
rational spatial structure for \(\gamma,f\), refined by edge strips.
Split full-line edge components at a rational point if needed. Every
cell is pointed: lineality would give opposite nonzero reference
recession vectors with opposite affine slopes, both supposedly positive.
Define \[
 D=\{p\in\mathbb{R}:p r_z\leq s_r\text{ for every cell recession generator}
                                      \ r\text{ with }r_G=0\},           \tag{12}
\] where \(s_r\) is the affine obstacle slope. Bounded \(u\) has zero
recession slope, so these are the same constraints for \(\gamma\). They
are finite, have rational variable coefficients and real constants, and
give a closed interval with zero in its interior by strict growth. For
\(p\in D\), \[
                         F_p(q)=\inf_{z\in\mathbb{R}}(f(q,z)-pz)          \tag{13}
\] is finite and attained at a vertex of some spatial cell fibre. The
fibre is pointed. Its recession cone is generated by zero-graph cell
rays: on a bounded strip all graph recession is zero; on an end strip it
is nonnegative, so positive graph rays cannot combine to yield zero
graph recession. The objective is nonnegative on the fibre recession
cone, so (8) supplies a minimizing vertex. There are finitely many
cells. If \(p\notin D\), a violating rational ray based at a rational
cell point gives \(F_p(q)=-\infty\) at a rational graph point. The
interval \(D\) is bounded: at any fixed rational graph point, finitely
many pieces give \(f(q,z)\leq M|z|+B\); if \(|p|>M\), (13) is
\(-\infty\). Therefore \(D\) is compact.

Subdivide graph edges at the projections of all spatial vertices. These
projections are rational and finite. Between consecutive projections, a
fibre vertex lies on a spatial one-face: intersection with the slicing
hyperplane lowers its minimal-face dimension by at most one, so a fibre
vertex's minimal face has dimension at most one. Such a traversing face
has rational trace \(z(t)=z_i+t b_i\); its affine obstacle height has
real coefficients. Its feasibility throughout that interval is constant,
since changes occur only at projected face vertices. Consequently \[
 F_p(t)=\min_i(A_i(p)+tB_i(p)),                              \tag{14}
\] with finitely many indices independent of \(p\), and affine
\(A_i,B_i\) with rational \(p\)-gradients and real constants. At a graph
node, direct fibre slicing gives the same kind of finite affine minimum.
Formula (14) extends continuously and concavely to the finite endpoints:
finite candidate limits have objectives at least \(F_p(v)\); conversely,
fixing any finite \(z\) and using continuity of \(f\) gives
\(\limsup F_p(t)\leq f(v,z)-pz\), then take its infimum. Thus changes of
endpoint fibre type are retained.

After the last projected vertex on an end, take its primitive parameter
\(t\geq0\) from that node. Define \[
 \theta_e(p)=\min_iB_i(p)=\lim_{t\to\infty}F_p(t)/t.          \tag{15}
\] Every \(B_i-\theta_e\geq0\), so \(F_p(t)-t\theta_e\) is
nondecreasing. Hence the line with slope \(\theta_e\) starting below
\(F_p(0)\) lies below the entire tail obstacle. All tails and pure/mixed
recession constraints are included.

For fixed \(p\in D\), on the finite graph node set \(V\) impose the
incidence system \[
 a_v\leq F_p(v),\qquad
 L_v(a,p)=\sum_{e=vw\text{ bounded}}\frac{r_e}{\ell_e}(a_w-a_v)
              +\sum_{e\text{ tail at }v}r_e\theta_e(p)\geq0.             \tag{16}
\] The actual primitive lengths \(\ell_e\) are positive rational
numbers. Every bounded cycle edge appears in this system; no graph is
unfolded.

It is feasible. The global competitor \(\gamma-C\leq f\) is actual PC.
Its partial infimum is finite on \(D\) by (12)--(13), and the scalar
lemma applies. Replace it on every bounded interval by its endpoint
chord and on every tail by its endpoint line of slope \(\theta_e\).
Convexity makes these replacements dominate it. Concavity of (14) keeps
the chords below \(F_p\); (15) keeps the tail lines below it. At both
endpoints the new outgoing chord derivative is at least the old
derivative. Convexity and the tail asymptotic upper bound force all old
tail secants to be at most \(\theta_e\), so its outgoing derivative also
increases. This proves (16).

Coordinatewise maxima preserve feasibility: at a node chosen from one
vector all neighbour values only increase, so its \(L_v\) does not
decrease. Choose a feasible witness \(b\); maximize \(\sum a_v\) in the
compact box \(b_v\leq a_v\leq F_p(v)\). If any feasible vector exceeds a
maximizing vector in one coordinate, their coordinatewise maximum is
feasible in that box and increases the objective. There is therefore a
greatest feasible vector \(a(p)\). Its chord/tail interpolation
\(k_p(q)\) is continuous, obeys (10), lies below \(F_p\), and dominates
every scalar function (11) lying below \(F_p\), by the same chord/tail
replacements.

Each \(a_v\) is a finite concave PA function on \(D\) with rational
gradients. Expand all caps into their affine minima and each weighted
sum of tail minima into a minimum over all choices of pieces. This turns
(16) into finitely many linear inequalities in \((a,p)\) with rational
variable coefficients and real constants. Eliminate all \(a_w\),
\(w\ne v\), by the proved elimination procedure. For fixed \(p\), the
section is \((-\infty,a_v(p)]\), since uniformly lowering all node
values preserves the Laplacians. Thus all nontrivial bounds on \(a_v\)
are upper bounds; zero-coefficient inequalities only restrict \(D\).
Dividing positive rational coefficients yields \[
                       a_v(p)=\min_j(c_{v,j}+g_{v,j}p),\quad
                       g_{v,j}\in\mathbb{Q}.                 \tag{17}
\] This proves continuity through the boundary of \(D\), as well as
rational gradients. The same holds for \(\theta_e\).

Set \[
                            Q(q,z)=\sup_{p\in D}(k_p(q)+pz).             \tag{18}
\] It is finite by compactness and continuity and actual PC by the
scalar lemma and the supremum Jensen inequality. It is at most \(f\).
For any raw \(h\leq f\) at rational \((q,z)\), the rational slice
\(h(q,\cdot)\) has the continuous convex completion proved above. A
touching affine support of slope \(p\) gives \(h(q,z)=k_{h,p}(q)+pz\).
The partial infimum is finite there, hence throughout the component by
propagation. A violating recession ray would contradict \(h\leq f\), so
\(p\in D\). Greatest-assignment domination gives \(h(q,z)\leq Q(q,z)\).

Partition \(D\) at the finitely many breakpoints of (17) and (15); let
\(J\) be the finite set of chamber endpoints. They may be real. On any
graph interval, \(k_p(q)+pz\) is affine in \(p\) on each chamber;
therefore (18) is the maximum over \(p\in J\) of finitely many real
affine strip functions. These agree at graph nodes. Hence \(Q\) is
continuous. At any real carrier point choose its minimal face on a
rational structure; unless it is a rational zero-face, rational vertices
of arbitrarily small enclosing simplices in that face have the point as
a real barycentre. The same-cell PCC and continuity of \(Q\) extend
\(h\leq Q\) to that point, without raw continuity. Applying this to
\(h=\gamma-C\) gives \[
                         \gamma-C\leq Q\leq f\leq\gamma+C.              \tag{19}
\] Thus \(Q\) itself is a raw competitor and bounds all competitors,
proving \(P_\gamma(u)=Q\) exactly.

\textbf{Rational spatial domains.} A finite real affine maximum alone
does not justify rational walls. Fix \(p_0\in J\). Let \(A_v\) be the
supergradient interval of \(a_v\) at \(p_0\), \(B_e\) that of
\(\theta_e\), and \(N_D\) the outward normal cone of \(D\) there.
Endpoints of \(A_v,B_e\) are rational by (17), and \(N_D\) has rational
generators by (12). For finite concave polyhedral \(b\), \(p_0\)
maximizes \(b(p)+pz\) on \(D\) exactly when \[
                                z\in-\partial^+b(p_0)+N_D.             \tag{20}
\] Sufficiency is the supergradient/normal inequality. If necessity
failed, separating the compact interval \(z+\partial^+b\) from the
closed cone \(N_D\) gives a direction nonpositive on all active normals
and strictly positive on every objective directional term. A small
feasible step raises the objective, a contradiction. This proves the
boundary cases too.

On a bounded edge \(0\leq t\leq\ell\), its full maximization domain is
\[
 z\in-(1-t/\ell)A_v-(t/\ell)A_w+N_D.                       \tag{21}
\] It is the convex hull of rational points
\((0,-\operatorname{Vert}A_v)\), \((\ell,-\operatorname{Vert}A_w)\),
plus the cone of rational vertical normal generators. On a tail the
domain is \[
 t\geq0,\qquad z\in-A_v-tB_e+N_D,                         \tag{22}
\] generated by rational base points and rational rays
\((1,-\operatorname{Vert}B_e)\) and vertical normal rays. Expanding the
finite weighted minima proves these supergradient formulas even at zero
coefficients. The domains are finite rational spatial polyhedra. Maxima
domains meet in equality faces; including their faces and removing
duplicates gives a rational complex. At a graph node every incident
formula has the same fibre domain \(-A_v+N_D\), so complexes glue.
Generated bounded-strip domains have vertices only at graph endpoints,
and tail domains only at their base. Thus \textbf{every} possible
zero-cycle vertex lies over a node of \(V\); there are no omitted
open-edge or distant-tail vertices.

By the PC-to-PAPC argument following (7), \(P\) is PAPC against all
local real cycles. The full recursive measure is finite and nonnegative.
The sequence \(P+1/n\) proves actual regularizability with exactly this
measure.

\textbf{Complete contact.} At any possible atom \(x=(v,z_0)\), suppose
\(P(x)<f(x)\). Define the nonempty compact exposed interval \[
 C_x=\{p\in D:a_v(p)+pz_0=P(v,z_0)\}.                      \tag{23}
\] It is convex because \(a_v\) is concave and the displayed expression
never exceeds its maximum. On \(C_x\), \(a_v=P(x)-pz_0\) is affine. Put
\[
 s_e(p)=(a_w(p)-a_v(p))/\ell_e\quad\text{on bounded outgoing edges},
 \qquad s_e(p)=\theta_e(p)\quad\text{on tails}.              \tag{24}
\] These are continuous concave polyhedral functions on \(C_x\), and
\(\sum r_es_e\geq0\). In local outgoing coordinates \(t\geq0\),
\(\zeta=z-z_0\), the homogeneous germ is \[
                            h_e(t,\zeta)=\max_{p\in C_x}(p\zeta+ts_e(p)).\tag{25}
\] Indeed the finite set \(J\) expresses all strip functions. Inactive
planes at \(x\) have positive gaps and stay inactive locally. Every
\(p\in C_x\) is a convex combination of endpoints of a chamber on which
all terms are affine; equality at the maximum forces its relevant
endpoints to be active. This proves the exact germ identity.

If \(C_x\) is a single point \(p_0\), after subtracting \(p_0\zeta\) the
germ is independent of \(\zeta\) on every branch. Its first divisor is a
constant-weight vertical line over the graph vertex (all interior branch
planes are affine); the function restricts to an affine function on that
line. Formula (5) gives zero second coefficient, or equivalently its two
opposite primitive vertical rays in (7) cancel. This retains artificial
cuts and all first faces by refinement compatibility.

If \(C_x\) is a nondegenerate interval, for every interior \(p\) choose
\(\eta>0\) with \([p-\eta,p+\eta]\subset C_x\). Globally in the vertex
fibre, \[
 P(v,y)-py\geq P(v,z_0)-pz_0+\eta|y-z_0|.                 \tag{26}
\] Thus continuous \(f(v,y)-py\) is coercive and attains its minimum. If
\(a_v(p)=F_p(v)\), (26) forces its minimizing point to be \(z_0\) and
equality \(f(x)=P(x)\), a contradiction. Hence \(a_v(p)<F_p(v)\)
throughout the interval interior. Greatest-assignment maximality forces
\(L_v(a(p),p)=0\): if positive, increase only \(a_v\) by a small amount
preserving its strict cap and that Laplacian. Every neighbour Laplacian
increases, tails are unchanged, and feasibility persists, contradicting
maximality. If there are no bounded incident edges, such an increase
changes no Laplacian and already contradicts the strict cap. Continuity
now gives \(\sum r_es_e=0\) on all of \(C_x\).

Positive weights and concavity force every \(s_e\) affine on this
interval: their nonnegative Jensen deficits have weighted sum zero. Its
affine gradient is rational, because an interior subinterval avoids the
finitely many walls of (17) and (15), whose gradients and primitive
lengths are rational. Write \(s_e(p)=\alpha_e+\beta_ep\),
\(\beta_e\in\mathbb{Q}\), \(\alpha_e\in\mathbb{R}\). The zero sum gives
\(\sum r_e\alpha_e=\sum r_e\beta_e=0\). Since the star direction kernel
is exactly this one relation, there exist a real linear form \(A\) and a
rational linear form \(B\) on its actual rational span, with
\(A(v_e)=\alpha_e\), \(B(v_e)=\beta_e\). This follows by factoring the
prescribed values through the rational direction matrix, not by
substituting a ray-generated lattice.

The full star germ is therefore \[
                              A(s)+\operatorname{supp}_{C_x}(\zeta+B(s)).\tag{27}
\] Extend these forms to the ambient space, rationally for \(B\). After
removing \(A\) and an affine multiple of the rational form
\(\ell=\zeta+B(s)\), the residual is a one-variable convex homogeneous
function, affine on \(\ell\geq0\) and \(\ell\leq0\), and zero on
\(\ell=0\). Refine the actual fan by this rational hyperplane. Its first
divisor vanishes away from \(\ell=0\), since the residual is ambient
affine there. On the entire support of the first divisor the residual is
zero, so its second divisor is zero by (5). This is the complete
recursion, including singular axes and all indices; it is not a
top-face-interior comparison. Consequently every noncontact atom has
coefficient zero. Finite support and bounded gap (19) prove (2).

For finitely many components construct the data, envelopes and measures
separately. A PCC's cells all contain its common face, so lie in one
connected component; a local cycle can similarly be treated
componentwise. Components are disjoint finite closed polyhedral subsets
and are open relative to their union. The finite component constructions
give a finite rational structure, a maximum of their finitely many
growth constants, and measures whose contact integrals add. Restore the
actual positive real product scalar on each component; zero stays zero.
This finishes the positive direction for all stipulated data.

\subsection{Negative direction: rational tripod projections on the
original
stars}\label{negative-direction-rational-tripod-projections-on-the-original-stars}

Assume both factors have branched junctions, choose \(A,B\), and
translate them to zero for local formulas. By the kernel proof the local
directions have sole primitive positive integer relations \(\lambda_i\),
\(\mu_j\), and actual weights \(a\lambda_i\), \(b\mu_j\). Partition each
incident index set into three nonempty groups \(I_k\), \(J_k\),
\(k=1,2,3\). Put \[
 S_k=\sum_{i\in I_k}\lambda_i,\quad T_k=\sum_{j\in J_k}\mu_j,
 \quad\nu_1=(1,0),\ \nu_2=(0,1),\ \nu_3=(-1,-1),\quad c_k=1,1,2.
\] Prescribe \[
 L_1(v_i)=\nu_k/S_k\ (i\in I_k),\qquad
 L_2(w_j)=\nu_k/T_k\ (j\in J_k).                            \tag{28}
\] These assignments annihilate the entire relation kernel, since their
weighted sums are \(\sum_k\nu_k=0\). Thus they factor to rational linear
maps on the actual rational spans, and extend to rational ambient maps
by a rational basis extension. Neither a lattice isomorphism nor an
integral map is claimed or required: the maps define a function only.
They do not replace the original cycles or saturated lattices.

Use ambient lattice coordinates to set, on the whole \(X\), \[
 H(p,q)=\|L_1(p-A)-L_2(q-B)\|_1,\qquad
 \rho(p,q)=\|p-A\|_1+\|q-B\|_1,
 \qquad\gamma=\max\{H,\rho-1\}.                            \tag{29}
\] A different positive rational cutoff scale is unnecessary. All affine
pieces have rational coefficients and all maximum walls are rational.
Intersect their dominance domains with the entire original product
structure, splitting full-line edge components at rational points if
present. Its cells are pointed and finite rational; denote this
reference structure by \(\Pi\). The function is continuous and ambient
convex. For every admitted PCC each affine piece has equality at the
barycentre, and taking its maximum gives Jensen. Hence \(\gamma\) is
actual PC on all structures, with real centres and remote endpoints. The
preceding PC-to-PAPC proof establishes all local real-cycle tests.
Alternatively, at any tested facet subtract an active supporting affine
piece; its tangent agrees with \(\gamma\), and each inward residual
derivative is nonnegative in (5), for arbitrary nonnegative real
weights.

Write \(\rho_0(r)=\|r_1\|_1+\|r_2\|_1\) for the homogeneous
ambient norm of a recession vector. For a nonzero recession vector
\(r\) of any reference affine cell, \[
 \gamma(z+tr)\geq\rho(z+tr)-1\geq t\rho_0(r)-\rho(z)-1.
\] Dividing by \(t\) gives \(\gamma_{\mathrm{lin}}(r)\geq\rho_0(r)>0\).
This proves the exact positive-growth condition on every pure and mixed
end, independent of the kernels of the projection maps.

There is rational \(r_0\in(0,1)\) such that \(\rho<r_0\) on \(X\) is
the intersection of the local product of original \(A,B\) stars with that neighbourhood, before any finite
other endpoint, and \(\gamma=H\) there. Indeed the original structure is
finite and closed; cells not incident to the chosen vertices have
positive distance from them, and the finitely many incident bounded
endpoints have positive distances. Choose a rational radius below these
distances and below one. Nonincident components cause no exception. In
this neighbourhood the active carrier and \(H\)-domain equations are
homogeneous about \((A,B)\). At every nonzero point, a sufficiently
small radial expansion and contraction stays in its minimal local face;
all far endpoint and cutoff inequalities remain strict. Such a point
cannot be a vertex of \(\Pi\). Therefore all \(\Pi\)-vertices other than
\((A,B)\) have \(\rho\geq r_0\).

Choose small rational \(\tau>0\) so that, for \(i\in I_k,j\in J_k\), \[
 z_{ij}=(A+\tau S_kv_i,\ B+\tau T_kw_j),\qquad
 \alpha_{ij}=\frac{\lambda_i\mu_j}{3S_kT_k}                 \tag{30}
\] are in their original incident product cells and
\(\rho(z_{ij})<r_0\). There are finitely many such positive distances,
so this is possible. Their coefficients sum to one (each group
contributes \(1/3\)); their two weighted displacements are
\(\tau\sum_i\lambda_iv_i/3=0\), \(\tau\sum_j\mu_jw_j/3=0\). Thus (30) is
a genuine rational-endpoint original-product PCC centred at \((A,B)\),
with the original product vertex as common face. Both projections equal
\(\tau\nu_k\), so \(H(z_{ij})=\gamma(z_{ij})=0\).

Choose an integer \(M\) with \(1/M<r_0\) and \(M\rho(z_{ij})\geq1\) for
every endpoint; define \[
                          u=(1-M\rho)_+.                  \tag{31}
\] It is bounded between zero and one and continuous real PA on finite
rational spatial domains. It vanishes at all endpoints (30), at every
other \(\Pi\)-vertex, and off this small star product. At \((A,B)\) its
value is one.

\subsection{Exact global raw equality, retaining cycles and all
ends}\label{exact-global-raw-equality-retaining-cycles-and-all-ends}

Every raw competitor \(q\leq\gamma+u\) obeys
\(q(A,B)\leq0=\gamma(A,B)\), by (30) and its zero endpoint obstacles. At
every other \(\Pi\)-vertex, \(q\leq\gamma\) by (31). To extend this to
any arbitrary real point \(x\) of any reference cell \(\sigma\), use
(8): \[
 x=\sum_i\beta_iv_i+\sum_jt_jr_j,\qquad
 \beta_i,t_j\geq0,\quad\sum_i\beta_i=1.
\] Choose \(\beta_{i_0}>0\), put \(T=\sum_jt_j\), and take integers
\(L\to\infty\) with \(T/L\leq\beta_{i_0}\). Replace that amount of
vertex coefficient by coefficients \(t_j/L\) on rational far endpoints
\(v_{i_0}+Lr_j\). This is a same-cell PCC centred exactly at \(x\), with
allowed real coefficients. The vertex bounds, the \textbf{sole upper}
growth \(q\leq\gamma+C_q\) at far endpoints, and affinity of \(\gamma\)
on \(\sigma\) yield \[
                         q(x)\leq\gamma(x)+(T/L)\max\{C_q,0\}.
\] Let \(L\to\infty\). Hence \(q\leq\gamma\) everywhere. Conversely
\(\gamma\) is a global actual PC competitor below \(\gamma+u\), with its
own upper growth bound. This proves \(P_\gamma(u)=\gamma\) on all
original bounded edges, cycles, ends and components. No raw continuity,
local competitor gluing, boundary omission or conical substitute is
used. The explicit \(\gamma+1/n\) approximation gives genuine PAPC
regularizability and exactly its PA operator.

\subsection{Complete primitive-normal junction
computation}\label{complete-primitive-normal-junction-computation}

In the original product quadrant \((s v_i,t w_j)\), \(s,t\geq0\), put
\(q_i=1/S_k\) if \(i\in I_k\), \(p_j=1/T_l\) if \(j\in J_l\). Directly
from the three specified tripod directions, \[
 H=c_kq_i s+c_lp_j t\quad(k\ne l),\qquad
 H=c_k|q_i s-p_jt|\quad(k=l).                              \tag{32}
\] Let \(a_i=(v_i,0)\), \(b_j=(0,w_j)\). The actual saturated plane
lattice is \(\mathbb Za_i\oplus\mathbb Zb_j\): a vector in this plane
belongs to the direct product ambient lattice exactly when its two
primitive factor-line coefficients are integral. Factor star spans may
be nonsaturated in their ray generators; this does not alter these
actual plane lattices.

Off-group quadrants are affine and have index-one axis cones. For a
same-group pair take a positive integer denominator \(D\), \(Q=Dq_i\),
\(P=Dp_j\), \(g=\gcd(Q,P)\). Its zero kink ray has actual primitive
vector \[
                            d_{ij}=(P/g)a_i+(Q/g)b_j.        \tag{33}
\] The integer coefficients are coprime; primitiveness follows also in
the whole ambient lattice because both factor vectors are primitive. The
two actual saturated cone indices are \[
                    \chi_{a_i,d_{ij}}=Q/g,\qquad
                    \chi_{b_j,d_{ij}}=P/g.                 \tag{34}
\] All top cone weights are \(\Omega_{ij}=ab\lambda_i\mu_j\), inherited
from the original graph product. At a kink, (6) and \(H(d_{ij})=0\) give
the full first weight \[
 D_H(d_{ij})=\Omega_{ij}\left(\frac gQ c_kq_i+\frac gP c_kp_j\right)
                             =2c_k\Omega_{ij}g/D.           \tag{35}
\] No denominator or index has been suppressed.

At an axis \(a_i\), the off-group neighbours are \(b_j\). For same-group
neighbours the representative quotient normal is
\(d_{ij}/(Q/g)=b_j+(p_j/q_i)a_i\). Therefore top balancing gives \[
 \sum\frac{\Omega_{ij}}{\chi}r_{ij}
       =ab\lambda_i\left(\sum_j\mu_jb_j+
                    \sum_{j\in J_k}\mu_j(p_j/q_i)a_i\right)
       =ab\lambda_i S_k a_i.                               \tag{36}
\] Since \(H(a_i)=c_k/S_k\), the off-group first term is
\(ab\lambda_i(4-c_k)\), and the tangential correction in (36) subtracts
\(ab\lambda_i c_k\). Thus \[
 D_H(a_i)=ab\lambda_i(4-2c_k),\qquad
 D_H(b_j)=ab\mu_j(4-2c_l).                                 \tag{37}
\] These and (35) are all rays in the local function fan and all are
nonnegative. Its balancing can be checked directly: kink contributions
in the first projection for a fixed \(i\in I_k\) sum to
\(2c_kab\lambda_i v_i\), adding its axis contribution gives
\(4ab\lambda_i v_i\). The second projection similarly gives
\(4ab\mu_jw_j\). Both total sums vanish by the original relations. This
check retains every same-group pair even when multiple original rays
share a projected direction.

Apply the second divisor (7). Kink rays have zero value. The entire
first-axis contribution is \[
 ab\sum_k(4-2c_k)c_k\sum_{i\in I_k}\lambda_i/S_k=4ab,
\] and the second-axis contribution is \(4ab\). Hence the
\textbf{complete} coefficient at the chosen junction is \[
                          (\gamma^2\cdot[X])\{(A,B)\}=8ab.              \tag{38}
\] This is the two-stage computation at the original singular product
junction, not a pushed-forward tripod calculation or a one apparent atom
claim.

\subsection{Full strictly negative integral and classification
conclusion}\label{full-strictly-negative-integral-and-classification-conclusion}

The complete PA measure is a finite zero-cycle on \(\Pi^{(0)}\). At
every other vertex the gap \(P-\gamma-u\) is zero; at \((A,B)\) it is
\(-1\). Formula (38) therefore gives the \textbf{full} integral
\(-8ab\). Its value requires no assumption about the masses of the other
contact vertices. Refining for the obstacle can insert vertices with
nonzero gap inside the small support; by successive refinement
compatibility their second coefficients are zero because they were not
vertices of \(\Pi\). Original lower faces, diagonal kink faces, bounded
cycles, all ends and other components are included in the two divisors
and global support argument. The finite PA zero-cycle has no added atom
at infinity; its constant-shift regularization has constant measures, so
no escaping or compactification boundary term occurs.

This proves failure of the universal property when both factors branch,
while the reconstructed graph-cylinder theorem proves it when either
factor is unbranched. The reference in the obstruction is flat on the
small balanced PCC curve. Consequently no prospective theorem requiring
an additional strict-convexity package is refuted. The original
fixed-reference arbitrary-balanced-pair P05.2 classification, all
individual references on branched products, higher-dimensional contact
packages, literal coefficients and analytic realizations remain
unresolved here.

\subsection{Evidence, reproducibility and failed
routes}\label{evidence-reproducibility-and-failed-routes}

\texttt{local-recursion.py} reads \texttt{local-input.json} and uses
exact standard-library rational arithmetic to enumerate every cone and
ray in two specified higher-valence local examples, including a factor
ray-generated lattice of index two. It calculates the actual indices in
(34), all top-balancing equations, every first coefficient, first-cycle
balancing and (7). The diagnostic divides out the formal scale \(ab\)
only in its input weights and reports normalized mass eight.
Substituting arbitrary actual scales multiplies every coefficient and
integral by \(ab\), as the proof shows. This finite computation checks
these local inputs; it does not certify the universal theorem or replace
the global proof.

The earlier native notes in the round-10 and round-14 campaign
\texttt{run.json} files were extracted from
\texttt{campaign.result.notes}, read in full, and preserved in
\texttt{sources/}. They are fallible comparison evidence, not supports
whose mathematical results are imported. The root localization proposal
and both predecessor arguments were read. All load-bearing graph,
coface, rationality, raw-envelope and operator claims are reconstructed
above. No numerical LP, total-mass agreement, old audit, receipt, or
granted theorem is used.

Failed or rejected routes retained in \texttt{research-notes.md}:
ray-generated lattices would lose (34); arbitrary real weight
coefficients in (16) without normalization would not justify (17); real
affine maxima without (20)--(22) would not justify rational spatial
walls; a small local atom without global support would not compute (3);
a graph-only tree lift would lose cycles; and the earlier coarse
rectangle roof fails a genuine PCC and is not a contact counterexample.
No new numerical search is promoted to proof.

For completeness, that last failure is reconstructed here rather than
used from a predecessor note. On the unit tripod with directions
\(\nu_k\) times \([0,1]_t\), restrict the ambient affine functions
\(A=1-t+(x-y)/2\) and \(B=t-(x-y)/2\) to the three unit radial
rectangles \(0\leq q\leq1\). The lower convex roofs of their cornerwise
maximum caps are \[
 R_1=\max\{1-q/2,1+q/2-t\},\quad
 R_2=\max\{1-q/2,q/2+t\},\quad R_3=1.
\] Each displayed plane lies below the four corner caps and meets three
corners; its two triangular domains (split by \(t=q\) for the first
rectangle and \(t=1-q\) for the second) cover the rectangle. Jensen on
those triangles bounds every convex corner minorant by the displayed
roof, proving the formulas. At the common centre \((0,1/2)\) all roofs
are one. The rational endpoints \((\nu_k/4,1/2)\), with weights \(1/3\),
are a genuine coface PCC with that centre; their roof values are
\(7/8,7/8,1\), with mean \(11/12<1\). Thus this coarse glued roof is not
PC, even though \(A,B\) and their functional maximum are PC by ambient
affine equality and the maximum Jensen inequality. Any corner assignment
dominating both arrays under their cornerwise maximum caps must equal
those caps, so this coarse roof-feasible set has no greatest assignment.
The true functional maximum has trace \(\max(1-t,t)\), adding the
missing fibre breakpoint. This proves a failed reduction, not regularity
or contact failure, and is not a premise for (2) or (3).

The complete frozen argument, finite diagnostic packet and independent
review receipt accompany this paper. The Xean native Codex review used
\texttt{gpt-6.1-sol} with \texttt{xhigh} reasoning and returned PASS
for the exact universal graph-product characterization, with no essential
external theorem premises. Its full report and the immutable
argument/source bindings were read and checked. The two wording
clarifications concerning the local star neighbourhood and homogeneous
recession norm are recorded; they do not change the mathematics.

The root read the complete proof, both canonical native notes and the
author audit. The two finite local examples were recomputed in a new
directory, preserving all original artifacts. All forty-nine saturated
cone indices were also checked directly by greatest common divisors of
the ambient two-by-two minors. These finite checks are supplementary
evidence; the universal theorem is the proof above.

The saved standalone source was opened in the native editor for
compilation. The exported PDF uses that same source. Computational
review, source bindings and successful typesetting are not formal
verification or priority certification. Original P05.2's classification
of fixed arbitrary balanced pairs remains open.

Bibliography: BMPS, arXiv:2603.08664v1 (2026), Definitions 2.1--2.4,
2.18--2.20, 2.28--2.31, 3.1--3.3, 3.14, 3.23, 4.4, 4.15, 6.2 and 6.9,
used to fix conventions only. The proof bodies inspected, and the
controlling book/predecessor snapshots, are indexed in
\texttt{source-boundaries.json}.

\end{document}
